% ============================================================================
\section{Route Planning}
\label{sec:route}
% Source: assignments/03_route/index.html
% ============================================================================

% ----------------------------------------------------------------------------
\subsection{Grid city}

Infinite grid of integer locations $(x,y)$; $s_{\text{start}}=(0,0)$;
$\text{Actions}(s)=\{(+1,0),(-1,0),(0,+1),(0,-1)\}$; $\text{Succ}(s,a)=s+a$;
$\text{Cost}((x,y),a)=1+\max(x,0)$; $\text{IsEnd}(s)=\indic{s=(m,n)}$ with $m,n\ge0$.

\problempart{1a}{Minimum cost}
\question{What is the minimum cost of reaching $(m,n)$ from $(0,0)$? Describe a path achieving it.
Is it unique?}
\expect{An expression for the minimum cost; 1--2 sentences describing one optimal path and whether it is unique.}
\begin{solution}
\todo{answer}
\end{solution}

\problempart{1b}{UCS on this problem}
\question{True or false:
(i) UCS never terminates because there are infinitely many states.
(ii) UCS returns the minimum-cost path and explores only $(x,y)$ with $0\le x\le m$, $0\le y\le n$.
(iii) UCS returns the minimum-cost path and explores only locations whose past cost is $\le$ the minimum cost to $(m,n)$.}
\expect{T/F for each sub-part.}
\begin{solution}
(i) \todo{} \quad (ii) \todo{} \quad (iii) \todo{}
\end{solution}

\problempart{1c}{UCS on an arbitrary graph}
\question{True or false:
(i) Adding a connection between two nodes cannot increase the minimum cost between any start and end node.
(ii) Making an action's cost small enough (possibly negative) guarantees that it appears in the path UCS returns, if a path exists.
(iii) Increasing every action's cost by 1 does not change the minimum-cost path, only its cost.}
\expect{T/F for each sub-part.}
\begin{solution}
(i) \todo{} \quad (ii) \todo{} \quad (iii) \todo{}
\end{solution}

% ----------------------------------------------------------------------------
\subsection{Finding shortest paths}

\problempart{2a}{\texttt{ShortestPathProblem}}
\codepart{3.2a}

\problempart{2b}{Stanford shortest-path experiments}
\question{Choose start locations and end tags with \texttt{get\_stanford\_shortest\_path\_problem}, and
visualise them. Give (i) a setting that gave new insight into travelling around campus (was the system useful?) and
(ii) a setting where the minimum-cost path is \emph{not} desirable. Is that caused by modelling assumptions
(missing tags or paths, inaccurate locations or distances)?}
\expect{Screenshots of both visualisations; one or more sentences on something learned about travelling, and
something incorrect about the map or the modelling assumptions.}
\begin{solution}
\todo{screenshots: figures/03\_route/2b\_*.png, plus discussion}
% \begin{figure}[H]\centering
%   \includegraphics[width=0.48\linewidth]{03_route/2b_insight.png}\hfill
%   \includegraphics[width=0.48\linewidth]{03_route/2b_undesirable.png}
%   \caption{...}\end{figure}
\end{solution}

\problempart{2c}{Negative externalities}
\question{What negative externalities could follow from wide deployment of this shortest-path system? Discuss
the effects on (1) users and (2) non-users, and ways to reduce the mismatch between the model and the real world
(see ``Why Traffic Apps Make Congestion Worse'', ``The Perfect Selfishness of Mapping Apps'', and the
Externalities module).}
\expect{3--5 sentences: two externalities (one for users, one for non-users) and at least one mitigation.}
\begin{solution}
\todo{answer}
\end{solution}

% ----------------------------------------------------------------------------
\subsection{Finding shortest paths with unordered waypoints}

\problempart{3a}{\texttt{WaypointsShortestPathProblem}}
\codepart{3.3a}

\problempart{3b}{Size of the state space}
\question{With $n$ locations and $k$ waypoint tags, what is the maximum number of states UCS could visit
under a suitable state definition for 3a?}
\expect{An expression in $n$ and $k$ with a brief justification.}
\begin{solution}
\todo{answer}
\end{solution}

\problempart{3c}{Stanford waypoints experiment}
\question{Choose a start location, waypoint tags (at least two) and an end tag with
\texttt{get\_stanford\_waypoints\_shortest\_path\_problem}, and visualise the route. What is unexpected or
interesting about it? Does it match your expectations, and what important factors does it fail to capture?}
\expect{A screenshot, the list of waypoint tags, and one or more sentences of discussion.}
\begin{solution}
\todo{screenshot figures/03\_route/3c\_*.png, tags used, discussion}
\end{solution}

\problempart{3d}{Ride-sharing and dual use}
\question{Ride-sharing apps use similar routing and have been criticised for exploitative nudges (e.g.\
``forward dispatch''). How could unordered waypoints improve the driver experience between a drop-off and the
next pick-up, and what driver information would you need? As a dual-use concern, how could a company use the
feature to raise profits at the expense of driver health, or what are the downsides of collecting that information?}
\expect{3--5 sentences: (a) waypoints that improve the driver experience; (b) at least one dimension of
\emph{driver identity} (not trip attributes such as time of day) that informs the choice; (c) a harmful use or a
downside of the data collection.}
\begin{solution}
\todo{answer}
\end{solution}

% ----------------------------------------------------------------------------
\subsection{Speeding up search with A*}
\codepart{3.4a--3.4c}
\texttt{a\_star\_reduction} (A* as UCS on modified costs), \texttt{StraightLineHeuristic},
\texttt{NoWaypointsHeuristic}.
